\hnheader[Neuronale Netze trainieren: Optimierer, Initialisierung, Normalisierung und Regularisierung.][Adam $\approx$ Momentum + adaptive Lernraten]{Deep Learning:}{Training \& Praxis}{Optimierer, Init, BatchNorm, Dropout, Early Stopping}
\hnsrc{main\_notes.pdf (Deep Learning: LayerNorm, SGD, Regularisierung); Adam, Dropout, BatchNorm, Initialisierung als Web-Ergänzung (markiert), Code: Vertiefungsseite}

\begin{hngrid}
%
\begin{hnp}[raster multicolumn=2,acc=hnorange]{1}{Momentum \& Adam (Web)}
\hnleg{$g$ Gradient, $\theta$ Parameter, $\alpha$ Lernrate, $v$ Geschwindigkeit, $m$ gleitender Mittelwert der Gradienten, $v$ (bei Adam) der Quadrate, $\beta,\beta_1,\beta_2$ Abklingraten, $t$ Schritt, $\varepsilon$ kleine Konstante.}
\textbf{Momentum}: $v\leftarrow\beta v+g$, $\theta\leftarrow\theta-\alpha v$ (glättet Zickzack).\\
\textbf{Adam}: $m\leftarrow\beta_1m+(1{-}\beta_1)g$, $v\leftarrow\beta_2v+(1{-}\beta_2)g^2$, $\hat m=\frac m{1-\beta_1^t}$, $\hat v=\frac v{1-\beta_2^t}$:
\[ \theta\leftarrow\theta-\alpha\frac{\hat m}{\sqrt{\hat v}+\varepsilon} \]
\end{hnp}
%
\begin{hnp}[raster multicolumn=2,acc=hnpurple]{2}{Lernraten-Schedules}
Konstante Lernrate ist selten optimal: \emph{Step decay}, \emph{Cosine}-Abklingen, \emph{Warmup} (klein starten, dann hoch). Zu große Rate $\to$ Verlust springt/divergiert; zu kleine $\to$ Stillstand. Vor dem Tuning: \hlt{Verlustkurve plotten}.
\end{hnp}
%
\begin{hnp}[raster multicolumn=2,acc=hnteal]{3}{Skizze: SGD vs.\ Momentum}
\centering
\begin{tikzpicture}[>=Stealth]
  \foreach \r in {.3,.6,.9}{\draw[gray!50] (2,1) ellipse ({\r*1.9} and {\r*.75});}
  \draw[hnred,thick] (.2,1.9)--(.6,.4)--(.9,1.6)--(1.3,.5)--(1.6,1.3)--(1.9,.9)--(2,1);
  \draw[hngreen,very thick] (.2,1.9) .. controls (1,1.5) and (1.5,1.1) .. (2,1);
  \node[font=\tiny,hnred] at (.55,.15) {SGD: Zickzack};
  \node[font=\tiny,hngreen] at (3.3,1.7) {Momentum: glatt};
\end{tikzpicture}
\end{hnp}
%
\begin{hnp}[raster multicolumn=3,acc=hnnavy]{4}{Initialisierung (Web)}
\emph{Nicht} alle Gewichte $=0$: dann sind alle Neuronen einer Schicht identisch (Symmetrie!). Zufällig mit passender Varianz ($n_{\mathrm{in}}$ Anzahl Eingänge): \textbf{Xavier} $\mathrm{Var}(W)=\frac{1}{n_{\mathrm{in}}}$ (Tanh/Sigmoid), \textbf{He} $\mathrm{Var}(W)=\frac{2}{n_{\mathrm{in}}}$ (ReLU).\\
Ziel: Aktivierungen und Gradienten weder verschwinden noch explodieren lassen.
\end{hnp}
%
\begin{hnp}[raster multicolumn=3,acc=hngreen]{5}{Normalisierung}
\textbf{BatchNorm} (Web): $\mathrm{BN}(z)=\gamma\frac{z-\mu_B}{\sqrt{\sigma_B^2+\varepsilon}}+\beta$ mit Batch-Statistiken $\mu_B,\sigma_B^2$ ($\gamma,\beta$ lernbare Skala und Verschiebung) (im Test: gleitende Mittelwerte).\\
\textbf{LayerNorm} (Notes): Statistik \emph{pro Beispiel} über die Features, unabhängig von der Batchgröße $\Rightarrow$ Standard in Transformern (Kapitel Neuronale Netze).\\
Wirkung: stabileres, schnelleres Training, größere Lernraten.
\end{hnp}
%
\begin{hnex}[raster multicolumn=3]{Beispiel: erster Adam-Schritt}
$g=2$, $\beta_1=0.9$, $\beta_2=0.999$, $\alpha=0.001$, $t=1$.\\
\hnsol\ $m=0.1\cdot2=0.2$, $v=0.001\cdot4=0.004$.\\
$\hat m=\frac{0.2}{0.1}=2$, $\hat v=\frac{0.004}{0.001}=4$.\\
Schritt $=\alpha\frac{2}{\sqrt4}=\ora{0.001}=\alpha$: Adam skaliert so, dass der erste Schritt etwa die Größe $\alpha$ hat, \emph{unabhängig} von der Gradientengröße.\\
\textbf{Dropout} ($p_{\mathrm{keep}}=0.5$): aus $(1,2,3,4)$ wird z.\,B.\ $(2,0,6,0)$ (verbleibende $\times2$, \emph{inverted dropout}).
\end{hnex}
%
\begin{hnp}[raster multicolumn=3,acc=hnrose]{6}{Regularisierung im Deep Learning}
\begin{itemize}
\item \textbf{Weight Decay} $=$ L2 (Notes: $\frac\lambda2\|\theta\|^2$).
\item \textbf{Dropout} (Web): im Training zufällig Neuronen ausschalten $\to$ implizites Ensemble.
\item \textbf{Early Stopping}: Training beenden, wenn der Validierungsfehler steigt.
\item \textbf{Datenaugmentierung}: Bilder drehen/spiegeln/zuschneiden.
\end{itemize}
\end{hnp}
%
\begin{hnp}[raster multicolumn=3,acc=hnpurple]{7}{Training debuggen}
\begin{enumerate}
\item Erst auf \emph{kleinem Batch} überanpassen (Verlust $\to0$?) -- sonst Bug.
\item Anfangsverlust prüfen ($\approx\log K$ bei $K$ Klassen).
\item Lernrate über Größenordnungen probieren.
\item Train- vs.\ Validierungskurve: Gap $\Rightarrow$ Overfitting.
\item Gradient Clipping bei explodierenden Gradienten.
\end{enumerate}
\end{hnp}
%
\begin{hnp}[raster multicolumn=3,acc=hnteal]{8}{PyTorch-Trainingsschleife (Skizze)}
\begin{pcode}[PyTorch]
\kw{for} Epoche \kw{in} Epochen:\\
\hii model.train()\\
\hii \kw{for} $(x,y)$ \kw{in} Mini-Batches:\\
\hiii opt.zero\_grad()\\
\hiii loss $=$ loss\_fn(model$(x)$, $y$)\\
\hiii loss.backward() \ \cm{\# Backprop}\\
\hiii opt.step() \ \ \ \ \cm{\# Update}
\end{pcode}
Vollständiger, getesteter Code: Vertiefungsseite.
\end{hnp}
%
\begin{hnp}[raster multicolumn=6,acc=hnorange]{9}{Key Takeaways}
\begin{hnchk}
\item Adam als robuster Standard, SGD+Momentum oft bessere Endleistung; Schedules und Warmup
\item He-Initialisierung für ReLU; BatchNorm/LayerNorm stabilisieren; Dropout, Weight Decay, Early Stopping gegen Overfitting
\end{hnchk}
\end{hnp}
%
\begin{hnp}[raster multicolumn=6,acc=hnpurple,colback=hnlilac!30]{?}{Selbsttest: kannst du das erklären?}
\hnquiz{Warum nicht alle Gewichte mit $0$ initialisieren?}{Alle Neuronen bleiben identisch (gleiche Gradienten): Symmetrie wird nie gebrochen.}{Was bewirkt Dropout?}{Zufälliges Ausschalten erzwingt redundante Merkmale, wirkt wie ein Ensemble.}{Was unterscheidet BatchNorm und LayerNorm?}{BatchNorm normiert über den Batch, LayerNorm über die Features eines Beispiels.}
\end{hnp}
%
\begin{hnbanner}[raster multicolumn=6]
„Training ist Handwerk: beobachte die Kurven, verändere eine Sache nach der anderen.“
\end{hnbanner}
\end{hngrid}
